\section{OLED display}

For quick access to information and data during testing, we used the \textbf{OLED
display} integrated in the ZedBoard. This display has a resolution of 128x32 pixels
and is controlled through I2C communication.

\section{Code implementation}

For the actual implementation of the OLED display, we used the following
repository:~\cite{OLED:repo} from Michael Mattioli, which provides a simple and
effective way to interface with the display implemented in VHDL and compatible with the ZedBoard.
\paragraph{}

Many of the functions provided by the repository were kept unchanged, while we focused
on \textbf{implementing new ones} to display the information we needed on the screen.
The OLED display is divided into 4 lines, each with a maximum of 16 characters,
and ASCII encoding is used to represent the characters.

\section{Displayed information}

Our audio processing module provides the \textbf{current frequency} of the audio signal. This data is
encoded in \textbf{Binary-Coded Decimal} (BCD) format, where each digit is written as 4-bits,
which guarantee the necessary range from 0 to 9.
The BCD format was chosen because it moves the operational weight of computing each digit
from the OLED writer code, to other modules already tasked with computations.
This allows a quick translation of the input data into an ascii encoded buffer, which will
then be used to print the current frequency.
\paragraph{}

Another information displayed is \textbf{the note} corresponding to the current frequency.
To achieve this we utilize the same data in BCD format and a \textbf{pre-compiled Look-Up Table} (LUT),
containing the various notes in a specific frequency range.
We chose the arbitrary value of 1000 Hz as the upper limit, which is more than enough to
represent any relevant note detected by our system.
\paragraph{}

The last information displayed on the OLED display is a simple \textbf{incremental counter},
used for debugging purposes, like ansuring the right refresh rate on the display.
This counter is also passed to the OLED module as a BCD encoded \texttt{std\_logic\_vector}.

